\chapter{LIMO : Robot Mobile}
\section{Introduction}

LIMO est un robot mobile innovant conçu pour répondre aux besoins éducatifs et de recherche dans le domaine de la robotique. Basé sur le framework ROS (Robot Operating System), il offre une plateforme polyvalente pour l'expérimentation et le développement de solutions autonomes.

\section{Caractéristiques Principales}

\subsection{Navigation Autonome et Perception}
LIMO intègre des capacités avancées de navigation autonome, incluant :
\begin{itemize}
    \item SLAM (Simultaneous Localization and Mapping) et V-SLAM (Visual SLAM)
    \item Détection d'obstacles
    \item Reconnaissance autonome de feux de signalisation
    \item Stationnement et manœuvres autonomes
\end{itemize}

\subsection{Architecture Matérielle Robuste}
\begin{itemize}
    \item \textbf{Structure métallique} pour une durabilité accrue
    \item \textbf{Moteurs puissants} intégrés aux roues
    \item \textbf{Espace d'expansion interne} de 0,9 L avec quatre ports USB série
\end{itemize}

\subsection{Quatre Modes de Direction Brevetés}
LIMO se distingue par sa capacité à adopter différentes configurations de direction :
\begin{itemize}
    \item Direction différentielle à quatre roues
    \item Direction à chenilles
    \item Direction omnidirectionnelle
    \item Direction de type Ackermann
\end{itemize}

\section{Spécifications Techniques}

\begin{table}[H]
\centering
\caption{Spécifications techniques du robot LIMO}
\begin{tabular}{|l|l|}
\hline
\textbf{Paramètre mécanique} & \textbf{Valeur} \\
\hline
Dimensions & 322 × 220 × 251 mm \\
\hline
Poids & 4,8 kg \\
\hline
Capacité de franchissement & 25° \\
\hline
\textbf{Système matériel} & \\
\hline
Interface d'alimentation & DC (5,5 × 2,1 mm) \\
\hline
Autonomie en fonctionnement & 40 min \\
\hline
Autonomie en veille & 2 h \\
\hline
\textbf{Capteurs} & \\
\hline
LIDAR & EAI X2L \\
\hline
Caméra & Orbbec DaBai \\
\hline
Ordinateur industriel & NVIDIA Jetson Nano (4 Go) \\
\hline
Module vocal & Assistant vocal IFLYTEK/Google Assistant \\
\hline
\textbf{Logiciel} & \\
\hline
Haut-parleurs & Double canal (2 × 2 W) \\
\hline
Écran & 7 pouces 1024 × 600 IPS tactile \\
\hline
Plateforme open source & ROS1/ROS2 \\
\hline
Protocole de communication & UART \\
\hline
\textbf{Équipements inclus} & \\
\hline
Roues & Tout-terrain × 4, Mecanum × 4, chenilles × 2 \\
\hline
\end{tabular}
\end{table}

\section{Environnement de Développement}

\subsection{Simulation}
LIMO est accompagné d'un environnement de simulation spécialement conçu, alimenté par Gazebo, permettant un développement et des tests facilités sans nécessiter le robot physique.

\subsection{Communauté et Ressources Open Source}
AgileX, la société derrière LIMO, maintient une communauté open source active offrant :
\begin{itemize}
    \item Packages ROS et ROS2 intégrés
    \item Démonstrations de programmation
    \item Manuel de développement robotique
    \item Spécifications produit détaillées
\end{itemize}

\section{Applications}

Grâce à sa polyvalence et ses multiples modes de direction, LIMO est adapté à diverses applications :
\begin{itemize}
    \item Recherche en navigation autonome
    \item Éducation à la robotique mobile
    \item Développement d'algorithmes de perception
    \item Prototypage de systèmes robotiques complexes
\end{itemize}

\section{Conclusion}

LIMO représente une plateforme complète et modulaire pour l'éducation et la recherche en robotique. Sa conception robuste, ses multiples modes de direction et son intégration native avec ROS en font un outil précieux pour explorer les défis de la robotique mobile autonome.
